\chapter{Reconstrucción de la dinámica cuántica mediante amplitudes coherentes sobre genealogías de ruta}

\epigraphtext{La onda no se añade a la historia: es la forma lineal en que la historia completa conserva simultáneamente todos sus caminos posibles.}{Principio de propagación HMT--MD}

La medida ya ha quedado definida como acoplamiento entre sistema, aparato y lector.
Ahora corresponde explicar qué se propaga antes de la publicación del resultado,
cómo una misma estructura sostiene onda y trayectoria persistente y de qué modo la
localidad observable convive con una genealogía global.

La respuesta de la Holografía Modular Triádica parte de una distinción que la
formulación puntual tiende a ocultar. Un estado no queda agotado por la casilla en
la que se encuentra ni por el número real que finalmente se lee. Conserva también
la ruta, la orientación, los cocientes, los acarreos, la frontera y el retorno que
lo llevaron hasta allí. Al linealizar ese espacio de historias, la suma de caminos
no aparece como una metáfora añadida a posteriori: es la expansión exacta del
propagador. La onda es la evaluación coherente de la totalidad de las rutas; la
partícula será, en el capítulo siguiente, la sección que persiste a través de una
de ellas.

\section{La amplitud global de \(1\) familia de caminos}

Sea \(X\) un conjunto finito de estados observables a una profundidad dada y sea
\(U:\mathbb C^X\to\mathbb C^X\) el operador de propagación de un paso. Una ruta de
longitud \(R\) desde \(i\) hasta \(f\) es una sucesión

\[
  \gamma=(x_0=i,x_1,\ldots,x_R=f)
\]

de transiciones admisibles. Su peso ordenado es

\[
  w_U(\gamma)=\prod_{t=1}^{R}U_{x_t x_{t-1}}.
\]

Aunque los pesos escalares conmutan, importa la secuencia indexada de transiciones:
no se está contando únicamente cuántas rutas alcanzan \(f\), sino qué factor asigna
cada paso a esa historia. La multiplicación compone las transiciones de una ruta; la
suma reúne historias alternativas que terminan en el mismo estado. Para pesos
operatoriales, el producto temporal es además genuinamente no conmutativo. Esta
diferencia entre producto dentro de una ruta y suma entre rutas es la versión lineal
de la doble lectura que atraviesa toda HMT.

\begin{theorem}[Suma finita de caminos]
Para todo \(R\geq 1\),
\[
  (U^R)_{fi}
  =
  \sum_{\substack{\gamma:i\to f\\|\gamma|=R}}
  w_U(\gamma).
\]
Si \(\psi_0\in\mathbb C^X\) es la amplitud inicial, entonces
\[
  \psi_R(f)
  =
  \sum_{i\in X}
  \sum_{\substack{\gamma:i\to f\\|\gamma|=R}}
  w_U(\gamma)\,\psi_0(i).
\]
\end{theorem}

\begin{proof}
La afirmación para \(R=1\) es la definición de \(U\). Al multiplicar una expansión
de \(U^R\) por \(U\), cada ruta de longitud \(R\) recibe exactamente una última
arista admisible. La distributividad suma todas las extensiones posibles y la
multiplicación conserva el orden de sus pesos. Por inducción, cada ruta de longitud
\(R+1\) aparece una sola vez. La segunda igualdad es la aplicación de la primera a
\(\psi_0\).
\end{proof}

Esta identidad tiene un contenido conceptual considerable. Una componente
\(\psi_R(f)\) no es una propiedad aislada del punto \(f\). Es la contracción de la
genealogía completa de rutas que lo alcanzan. Dos historias pueden acabar en el
mismo valor visible y, sin embargo, contribuir con fases opuestas; también pueden
reforzarse. La interferencia consiste precisamente en sumar amplitudes antes de
formar un observable positivo. No se necesita introducir una onda material
adicional que recorra un espacio distinto del de las historias: la estructura
ondulatoria está en la composición lineal de las rutas.

Cuando \(U\) es unitario,

\[
  \|\psi_R\|_2=\|U^R\psi_0\|_2=\|\psi_0\|_2.
\]

La cantidad total que después puede normalizarse como probabilidad permanece
constante. Si el lector físico asigna a una salida \(f\) el peso

\[
  p_R(f)=\frac{|\psi_R(f)|^2}{\|\psi_R\|_2^2},
\]

entonces \(p_R\) es una distribución normalizada y su evolución es compatible con
la conservación unitaria. Éste es el punto exacto en el que la amplitud de rutas se
convierte en estadística de resultados. El capítulo de la medida ya ha localizado
el acoplamiento que publica una salida; aquí queda identificado el objeto que ese
acoplamiento evalúa.

La construcción también aclara por qué un número visible no basta para reconstruir
la física. El mapa

\[
  \{\text{rutas profundas}\}\longrightarrow X
\]

olvida, en general, la fase relativa entre caminos. Conocer únicamente la
distribución final de casillas no permite invertir la suma ni recuperar la
genealogía. El estado HMT completo conserva lo que la pantalla contrae. Por eso la
memoria no es un adorno interpretativo: es el dato del que depende la interferencia.

En el límite de refinamiento, la pregunta correcta no es si la ruta discreta
``se vuelve'' misteriosamente continua, sino si los propagadores forman una familia
compatible. Para proyecciones de refinamiento \(p_{m+1,m}:X_{m+1}\to X_m\), la
coherencia exige que los observables cilíndricos y los estados transportados
concuerden. Cuando esa compatibilidad es exacta, la amplitud de una historia gruesa
es la suma de las amplitudes de todas sus prolongaciones finas. Cuando los defectos
son sumables, las expectativas de observables cilíndricos uniformemente acotados
forman una familia de Cauchy. Así se obtiene continuidad operacional sin borrar la
estructura discreta que la genera.

\begin{principle}[La onda como memoria lineal]
La función de onda HMT a profundidad finita es la evaluación lineal de la extensión
de rutas. Conserva de forma coherente las alternativas compatibles con la
preparación y la frontera; el valor observado es una proyección posterior de esa
amplitud global.
\end{principle}

Esta formulación es más fuerte que afirmar que ``todo ocurre por muchos caminos''.
Especifica el dominio, el peso de cada camino, la regla de composición y la norma
conservada. También contiene su control negativo: si una supuesta amplitud no
respeta la composición de rutas, cuenta dos veces una historia, consulta el
resultado futuro para asignar el peso o pierde la norma bajo un propagador
declarado unitario, no es la función de onda construida aquí.

\section{Onda piloto, ruta realizada y doble orientación}

La estructura anterior contiene ya las dos piezas que una teoría de onda piloto
necesita distinguir. La amplitud \(\psi_t\) organiza globalmente las alternativas;
una ruta persistente \(x_t\) registra la historia efectivamente realizada. La onda
y la partícula no son dos sustancias rivales. Son dos niveles de lectura del mismo
estado enriquecido: el campo lineal de posibilidades y la sección que conserva su
identidad al atravesarlo.

La ley de guía debe convertir esta arquitectura cinemática en dinámica. Su forma
mínima es un mapa

\[
  G:(\psi_t,x_t)\longmapsto x_{t+1},
\]

o, si la actualización es estocástica, un núcleo
\[
  P_{\psi_t}(y\mid x).
\]

La condición central no es verbal, sino matemática: la distribución guiada debe ser
equivariante con la evolución de la amplitud,

\[
  |\psi_{t+1}(y)|^2
  =
  \sum_x P_{\psi_t}(y\mid x)\,|\psi_t(x)|^2.
\]

Esta igualdad asegura que una población inicialmente distribuida según
\(|\psi_t|^2\) continúa distribuida de ese modo después de un paso. La guía no puede
elegirse para reproducir cada detector retrospectivamente; ha de quedar fijada por
el estado, el propagador, la orientación y el registro antes de conocer la salida.
HMT añade una posibilidad que una descripción puntual no posee: \(G\) puede
depender de cociclos de ruta y de memoria que no descienden al valor visible
\(x_t\). La acción local puede así estar informada por una genealogía global sin que
aparezca una fuerza suplementaria que viaje desde el futuro.

La lectura bidireccional profundiza esta estructura. Para una ruta orientada
\(\gamma\) con

\[
  \partial\gamma=x_{\rm abs}-x_{\rm em},
\]

la involución de retorno define

\[
  \gamma^\vee=C_{\rm rev}\gamma,
  \qquad
  \partial\gamma^\vee=x_{\rm em}-x_{\rm abs}.
\]

No se trata de dos caminos independientes elegidos libremente. Son las dos
orientaciones de una misma historia. La longitud interna es invariante,

\[
  L_{\rm int}(\gamma^\vee)=L_{\rm int}(\gamma),
  \qquad
  L_{\leftrightarrow}=2L_{\rm int}(\gamma),
\]

y la inversión intercambia los canales conjugados. La preparación y la absorción
dejan de ser extremos ontológicamente desconectados: pertenecen a una clausura en
la que la ida y el retorno deben ser compatibles.

Cuando esta involución se representa en la teoría física de propagadores, las dos
orientaciones admiten la lectura retardada y avanzada. La componente retardada
transporta la preparación hacia la absorción; la avanzada transporta la condición
de retorno hacia la emisión. La historia observable es su ajuste mutuo. Esta
formulación expresa de manera precisa la intuición de una onda piloto con adelanto
y retardo: la guía local recibe información de las dos fronteras del mismo objeto
orientado.

No hay en ello una autorización para enviar mensajes controlables hacia el pasado.
Una solución avanzada de una ecuación y una señalizable elección futura son objetos
distintos. La representación física debe demostrar que los propagadores satisfacen
la ecuación de evolución, que las condiciones de contorno se componen
consistentemente y que las probabilidades marginales no permiten modular una
salida anterior mediante una decisión posterior. La doble orientación es el
soporte; la causalidad observable es una propiedad que debe verificarse sobre su
lector físico.

La correlación entre ambas hojas puede formalizarse en dos espacios de Hilbert,
\(\mathcal H_{\rm adv}\) y \(\mathcal H_{\rm ret}\). Para una distribución
\(p=(p_i)\), el estado

\[
  |\Psi_p\rangle
  =
  \sum_i\sqrt{p_i}\,
  |i\rangle_{\rm adv}\otimes|i\rangle_{\rm ret}
\]

tiene matrices reducidas con espectro \((p_i)\) y entropía

\[
  S_{\rm ent}(p)=-\sum_i p_i\log p_i.
\]

La igualdad es exacta una vez declarada la factorización bicapa. Una única historia
posee rango de Schmidt uno y entropía nula; varias historias correlacionadas
producen entropía positiva. El resultado no identifica sin más esta entropía con
\(\log\phiHMT\), con la información de una celda triádica ni con la entropía
termodinámica. Cada una pertenece a un objeto diferente. Lo que demuestra es que
ida y retorno pueden formar un estado no separable y que la información compartida
se calcula sin metáforas.

\begin{test}[Guía bidireccional]
Una realización física de la onda piloto HMT debe fijar, antes del contraste, el
propagador \(U\), la ley \(G\) o el núcleo \(P_\psi\), la involución
\(C_{\rm rev}\) y las condiciones de frontera. Queda refutada si pierde
equivarianza, asigna dos rutas incompatibles al mismo estado completo, necesita
conocer el detector futuro para escoger la trayectoria o permite señalización
controlable contra la orientación causal publicada.
\end{test}

Esta prueba no relega la propuesta a una simple posibilidad. Al contrario: muestra
que el núcleo matemático ya está construido y señala el único lugar donde debe
actuar la dinámica física. La suma de caminos no está pendiente; la ruta persistente
no está pendiente; la inversión tampoco. Lo que se decide experimentalmente es qué
Hamiltoniano realiza la guía y qué condiciones de contorno convierten las dos hojas
en las ondas avanzada y retardada de la naturaleza.

\section{Bell: localidad del dato y globalidad de la genealogía}

La distinción entre estado completo y valor proyectado permite plantear con
precisión el problema de Bell. Considérense ajustes \(a,b\in\{0,1\}\), resultados
\(A,B\in\{-1,+1\}\) y un estado compartido \(\lambda\). Una teoría local
factorizable satisface

\[
  P(A,B\mid a,b)
  =
  \int_\Lambda
  P_A(A\mid a,\lambda)\,
  P_B(B\mid b,\lambda)\,
  \dd\rho(\lambda),
\]

con una medida independiente de los ajustes,
\[
  \dd\rho(\lambda\mid a,b)=\dd\rho(\lambda).
\]

Definiendo
\[
  E_{ab}
  =
  \sum_{A,B=\pm1}AB\,P(A,B\mid a,b),
\]

se obtiene
\[
  S_{\rm CHSH}=E_{00}+E_{01}+E_{10}-E_{11},
  \qquad |S_{\rm CHSH}|\leq2.
\]

HMT conserva este teorema. Añadir memoria oculta a \(\lambda\) no rompe la cota si
se mantienen la factorización, la independencia de ajustes y la ausencia de
postselección. Este control es importante: impide convertir la riqueza de TPK en
una explicación automática de cualquier correlación.

La contribución holográfica se formula de otro modo. Las dos alas de un experimento
pueden ser proyecciones locales de una misma sección enriquecida, con frontera y
genealogía comunes. El dato publicado en cada ala es local; el estado del que ambos
descienden no tiene por qué factorizarse en dos historias profundas independientes.
La no localidad no consiste entonces en una señal que cruza el espacio después de
elegidos los ajustes, sino en la imposibilidad de descomponer la preparación global
en dos estados completos autónomos.

La no señalización y la localidad de Bell no son equivalentes. Puede cumplirse

\[
  \sum_B P(A,B\mid a,b)=P(A\mid a),
  \qquad
  \sum_A P(A,B\mid a,b)=P(B\mid b),
\]

y, sin embargo, fallar la factorización anterior. Éste es el dominio exacto que una
realización HMT debe ocupar: correlación global sin canal de comunicación
superlumínico. El lector contextual puede depender del ajuste en la forma en que
proyecta el estado común, pero la distribución marginal de un observador no puede
depender de la elección remota.

\begin{construction}[Protocolo HMT para Bell]
Se fija primero un estado enriquecido \(x\), una medida de preparación y cuatro
lectores \(L_{ab}\). Después se declara un mapa binario de resultados y se calcula
la tabla completa
\[
  P_{\rm HMT}(A,B\mid a,b).
\]
La construcción es físicamente válida sólo si queda congelada antes de consultar
los datos, está normalizada, carece de postselección oculta y publica a la vez
\(S_{\rm CHSH}\) y los cuatro controles marginales de no señalización.
\end{construction}

El denominado Algoritmo 3 pertenece a esta construcción concreta. Su objetivo no es
volver a demostrar el teorema clásico, sino construir el descenso desde la
genealogía compartida hasta una distribución contextual. Si conserva todas las
premisas de Bell, producirá necesariamente \(|S|\leq2\). Si produce
\(|S|>2\), debe mostrar de forma explícita qué factorización no desciende y por qué
la no señalización sí se conserva. Ésta es la explicación holográfica completa del
problema de localidad: el espacio observable puede separarse aunque la historia que
lo genera no sea separable.

La propuesta es especialmente natural en una teoría donde suma y producto son
lectores diferentes. La factorización de Bell intenta expresar una probabilidad
conjunta como producto condicionado de respuestas locales; la amplitud cuántica
surge, en cambio, de sumar historias antes de proyectarlas. Si la proyección se
realiza demasiado pronto, se pierde la fase común y sólo queda el politopo local.
Si se conserva el estado global hasta la lectura conjunta, sobreviven términos de
interferencia que no admiten aquella descomposición. La incompatibilidad no es una
licencia para violar la lógica: es la diferencia tipada entre componer alternativas
por suma y componer probabilidades locales por producto.

\section{Del teorema de rutas al contraste cuántico}

La arquitectura cuántica obtenida se articula mediante una cadena única de mapas. HMT
construye un espacio discreto enriquecido de historias; su linealización produce la
suma coherente de caminos; la unitariedad conserva la norma; una sección
persistente permite hablar de trayectoria; la inversión orientada entrega el par de
ida y retorno; la factorización bicapa cuantifica su correlación; y el análisis de
Bell distingue la localidad del dato de la separabilidad de la genealogía. No son
analogías yuxtapuestas, sino mapas sucesivos sobre el mismo estado.

La teoría ofrece así una lectura unificada de los rasgos cuánticos. La
superposición es la coexistencia lineal de prolongaciones; la interferencia es su
suma con fase; la partícula es la persistencia de una sección; la aparente reducción
es la selección de una fibra por el lector; el entrelazamiento es la no
factorización de hojas o subsistemas; y la no localidad es la globalidad de la
genealogía conservada bajo proyecciones localmente no señalizantes.

Los contrastes físicos no se añaden para debilitar esta lectura, sino para decidir
su realización natural. Una red interferométrica finita puede implementar \(U\) y
comparar la expansión de caminos con el patrón final. Una reconstrucción de
corrientes puede someter a prueba la ley de guía. La inversión controlada de las
condiciones de frontera puede comprobar la reciprocidad
\(\gamma\leftrightarrow\gamma^\vee\). Un experimento Bell con lectores sellados
puede decidir si la memoria profunda genera una distribución no factorizable y no
señalizante.

\begin{test}[Certificado conjunto de compatibilidad cuántica]
El certificado distingue cuatro condiciones independientes de validez:
\begin{enumerate}
  \item el propagador no reproduce por suma coherente las amplitudes observadas;
  \item ninguna ley de guía natural es equivariante con \(|\psi|^2\);
  \item la doble orientación no admite propagadores físicos recíprocos sin
  señalización causal indebida;
  \item el descenso Bell es local factorizable o viola la no señalización.
\end{enumerate}
El fallo de una interfaz localiza el mapa incorrecto; no borra los teoremas finitos
de rutas que la preceden.
\end{test}

La tesis positiva es, por tanto, nítida. La cuántica no obliga a abandonar la
historia de un sistema; obliga a enriquecer qué se entiende por historia. Un punto
sin genealogía no puede portar por sí solo interferencia, retorno ni correlación.
Una extensión orientada sí puede hacerlo. Su onda es la memoria coherente de todos
los caminos compatibles; su partícula es el cierre que persiste. El tránsito entre
ambas formas no es un salto ontológico: es el paso de la totalidad lineal de rutas a
una sección realizada de esa misma totalidad.
